% Why a corrected clock is slewed and never stepped.
\begin{tikzpicture}[font=\sffamily\scriptsize, x=1cm, y=1cm]

% ================================================================ stepped
\node[signame, anchor=east] at (-0.3,3.4) {stepped};
\draw[taxis] (0,2.0) -- (7.4,2.0);
\draw[taxis] (0,2.0) -- (0,4.6);
\node[font=\sffamily\tiny, anchor=east, rotate=90] at (-0.35,3.3) {reported time};

% the clock runs, then jumps back
\draw[signal, line width=1pt] (0,2.2) -- (3.4,4.25);
\draw[signal, line width=1pt] (3.4,3.35) -- (7.2,4.35);
\draw[faultedge] (3.4,4.25) -- (3.4,3.4);
\node[tlabel, anchor=south, text=red!60!black] at (2.9,4.45) {offset applied as a jump};

% the duplicate-timestamp defect
\draw[tick] (0,3.62) -- (7.4,3.62);
\node[sample] at (2.35,3.62) {};
\node[sample] at (4.42,3.62) {};
\node[tlabel, anchor=south] at (2.35,3.68) {event 1};
\node[tlabel, anchor=south] at (4.42,3.68) {event 2};
\node[tlabel, anchor=west, text=red!60!black] at (5.2,3.3)
  {both are recorded at the same instant};
\draw[faultedge, <->] (2.35,3.5) -- (4.42,3.5);

% ================================================================ slewed
\node[signame, anchor=east] at (-0.3,0.6) {slewed};
\draw[taxis] (0,-0.8) -- (7.4,-0.8);
\draw[taxis] (0,-0.8) -- (0,1.8);
\node[font=\sffamily\tiny, anchor=east, rotate=90] at (-0.35,0.5) {reported time};

\draw[signal, line width=1pt] (0,-0.6) -- (3.4,1.45);
% a slightly steeper segment: the correction absorbed over a bounded window
\draw[signal, line width=1pt, draw=blue!55!black] (3.4,1.45) -- (5.6,1.62);
\draw[signal, line width=1pt] (5.6,1.62) -- (7.2,2.0);
\draw[tspan] (3.4,-0.45) -- node[tlabel]{the correction window, bounded} (5.6,-0.45);
\node[tlabel, anchor=west] at (5.75,1.35) {rate returns to normal};

\draw[tick] (0,0.95) -- (7.4,0.95);
\node[sample] at (2.35,0.95) {};
\node[sample] at (4.65,0.95) {};
\node[tlabel, anchor=south] at (2.35,1.01) {event 1};
\node[tlabel, anchor=south] at (4.65,1.01) {event 2};
\node[tlabel, anchor=west] at (5.3,0.62) {ordering preserved throughout};

% ================================================================ the readings
\node[umlnote, text width=50mm, anchor=north west] at (0,-1.4)
  {The stepped clock is one line of code and it is what most first attempts do.
   It makes time non-monotonic, so two events that happened at different moments
   carry the same timestamp. Both cited papers say the same thing about this: in
   a real-time system that situation must not be allowed, and the fix is to
   adjust the rate rather than the value.};
\node[umlnote, text width=46mm, anchor=north west] at (5.4,-1.4)
  {The slewed clock is knowingly wrong for the length of the window, and it is
   never inconsistent. The rate change is clamped, so a node told it is a minute
   out reports the discrepancy rather than spending a minute catching up: a
   large step almost always means the other machine is wrong.};
\end{tikzpicture}
